\section{Euler and Other PDEs as Departure Readings}

Euler is not a second Silver root and is never used by starting from Euler and
working backward.  Begin with the same smooth-data, fixed-viscosity
Navier--Stokes history and carry that object to the candidate departure.  The
only Silver question is whether the object has left its original full VPI
participation.

If the forward limit from Navier--Stokes actually removes the viscous term and
leaves no Reynolds stress, forcing, source, pressure, incompressibility,
compactness, or ancestry defect, then that same departure is governed by
Euler.  If one of those residuals survives, the honest limiting equation is
the corresponding Euler--Reynolds, Navier--Stokes--Reynolds, forced, or defect
law.  If the viscous coefficient remains positive and the entire same-history
VPI structure remains intact, no Euler departure has occurred.

These alternatives do not enlarge the logic of Silver.  Each successful
law-identification theorem says only how the original participation was lost.
Together with Dead, Blown, Jump, Part, and Field, they exhaust the negative
side
\[
  \neg\operatorname{Smooth}(Q)
  \Longrightarrow
  \neg\Member(Q)
  \Longrightarrow
  \neg\operatorname{VPIParticipation}(Q).
\]
The Silver conclusion is its contrapositive:
\[
  \operatorname{VPIParticipation}(Q)
  \Longrightarrow
  \Member(Q)
  \Longrightarrow
  \operatorname{Smooth}(Q).
\]

A weak record that merely satisfies Navier--Stokes distributionally is not by
that fact a full VPI participant: full participation retains the same material
history and coherent field realization as well as the equation.  Hence weak
fixed-viscosity solvability creates no additional nonsmooth-participant branch.
